% CUE_GCT isolation track. Run CUE-GCT-20261121-I.
% Model: GPT-6 Astra Pro; configured effort is not exposed to this agent.
% Bibliographic keys required: Hoza2026, AGKS2014, ForbesShpilka2013, MulmuleyGCTV, AllenderOgihara1996.

\section{Isolation: the exact transfer, the quotient, and the surviving gate}
\label{sec:isolation-track}

We use a field $F$, $n$ variables, individual degree at most $\delta$, and
an ROABP expression
\[
 f(\mathbf x)=u^{\mathsf T}A_1(x_{\pi(1)})\cdots A_n(x_{\pi(n)})v.
\]
The order $\pi$ may be unknown.  A polynomial-size hitting set here means
one set depending only on $(n,w,\delta,F)$, not on the matrices or their
coefficients, which hits every nonzero member of the class.  Arithmetic
complexity and bit complexity are distinguished throughout.

\paragraph{Background I-B01.}
The stipulated new graph theorem is Hoza (2026), Theorem 1.3, pp.~2--3;
its quantitative version is Theorem 2.4, p.~5.  For every simple layered
digraph of width $W$ and length $N$, some seed in a family of size
$(WN)^{O(1)}$ gives a minimum-weight path unique for every reachable pair.
The weights are nonnegative integers bounded by $B=(WN)^{O(1)}$ and are
computable in logarithmic space.  A path weight sums its internal vertex
weights.  With a binary code $C$ of length $\ell$ and distance $d$, the
formula is
\[
 \rho_x(a,i)=6^{\nu_2(i)}\sum_{j=1}^{\ell}6^{\ell-j}C(a)_j x_j
 \quad(i>0),\qquad \rho_x(a,0)=0,
\tag{I.1}\label{eq:I-hoza}
\]
and the probability of a bad uniform seed is at most
$2^{-d}W^4N^2$.  The graph is fixed before the seed is chosen.
We use this result as supplied; the present work does not claim an
independent verification of its complete theorem.

\paragraph{Background I-B02.}
Agrawal, Gurjar, Korwar, and Saxena (2014), \S3, Definition~5, call
$\omega\in\mathbb Z_{\geq0}^{n}$ basis isolating for
$D=\sum_a v_a\mathbf x^a$ if there is a coefficient basis indexed by $B$
whose monomial weights are distinct, and each $v_a$ with $a\notin B$
belongs to the span of basis coefficients of strictly smaller weight.
Their Lemma~6 preserves every nonzero scalar projection under
$x_i=t^{\omega_i}$; Lemma~8 and Theorem~1 give their quasipolynomial
construction, including unknown variable order.  Those known results
are cited, not claimed as findings here.

\subsection{I-01. A genuine polynomial transfer for vertex-variable path polynomials}
\label{I-01}

\paragraph{Statement.}
Let $G$ range over simple layered digraphs with fixed vertex slots,
width $W$, length $N\ge1$, fixed source $s$ and target $t$, and arbitrary
scalar edge labels $c_e\in F$.  Give every internal vertex slot its own
variable $z_v$, and form
\[
 P_G(\mathbf z)=\sum_{p:s\leadsto t}\left(\prod_{e\in p}c_e\right)
                        \prod_{v\in\operatorname{int}(p)}z_v.
\]
There is a uniform hitting set of cardinality at most
\[
 2^{\ell}\bigl((N-1)B+1\bigr)=(WN)^{O(1)}
\tag{I.2}
\]
over any field containing $(N-1)B+1$ distinct elements.  Its points have
coordinates $z_v=a^{\rho_x(v)}$, where $x$ runs over Hoza's seeds and
$a$ runs over a fixed set of $(N-1)B+1$ field elements.  The same conclusion
holds for independent edge variables after subdividing edges; the new
width is at most $W^2$ and the new length is $2N$.

\paragraph{Proof.}
Delete zero-labelled edges.  In a simple layered digraph two different
$s$--$t$ paths have different sequences of internal vertices and hence
different vertex monomials, except in the length-one case, where there
is at most one edge.  Thus $P_G\ne0$ exactly when some path survives.
Choose a min-isolating seed for the surviving graph using I-B01.  In
$P_G(t^{\rho_x(v)})$, the coefficient of the least path weight is the
product of labels on its unique minimum path, and is nonzero in the
field.  The degree is at most $(N-1)B$.  A nonzero polynomial of that
degree has at most $(N-1)B$ roots.  This proves (I.2), uniformly in the
edge labels.  For edge variables, insert one distinct intermediate
vertex on each edge and give it that edge's variable; give the original
internal vertices auxiliary independent variables as well.  The same
argument provides a hitting set for the expanded vertex-variable
polynomial.  To obtain points only on edge variables, assign to an edge
the product of its intermediate-vertex value and the value of its head
when the head is an internal original vertex.  Each original path then
receives exactly the product of all its expanded internal-vertex values.
This combines exponents additively, leaves a polynomial degree bound,
and proves the edge-variable statement. \hfill$\square$

\paragraph{Scope.}
This is an exact consequence of the new graph theorem.  It does not
identify several vertex or edge variables with the same ROABP variable.
That identification is the substantive next step, treated below.
Over a smaller finite field, the argument gives an extension-field
hitting set after constructing a sufficiently large extension; it is
not a statement that the same number of base-field points exists.

\subsection{I-02. The quotient obstruction is already present in a two-path diamond}
\label{I-02}

\paragraph{Statement.}
Let the monomial map $q:F[z_1,\ldots,z_s]\longrightarrow
F[x_1,\ldots,x_n]$ send $z_j$ to $\mathbf x^{a_j}$, and let $A$ be
the matrix with columns $a_j$.  A diagonal one-parameter action
$z_j\mapsto t^{r_j}z_j$ is induced by a diagonal action
$x_i\mapsto t^{\omega_i}x_i$ precisely when
$r=A^{\mathsf T}\omega$.  If the $a_j$ generate a proper sublattice,
requiring $r$ to annihilate $\ker(A:\mathbb Z^s\to\mathbb Z^n)$
only guarantees a character on the image lattice; an integral
extension to all $\mathbb Z^n$ is an additional lattice condition.
Consequently path isolation does not, in general, descend through $q$.

\paragraph{Proof and counterexample.}
Equivariance on each generator reads
$t^{r_j}\mathbf x^{a_j}=t^{\langle\omega,a_j\rangle}\mathbf x^{a_j}$,
which is equivalent to the displayed matrix equation.  The kernel
description follows from the factorization property of homomorphisms
of free abelian groups, with the stated extension qualification.

For the counterexample take the simple diamond
$s\to u\to t$ and $s\to v\to t$.  Label the two paths so that their
vertex-variable polynomial is $z_u-z_v$.  The map $q(z_u)=q(z_v)=x$
gives the zero polynomial.  Every min-isolating weight on this diamond
has $r_u\ne r_v$, so that its weighted polynomial
$t^{r_u}-t^{r_v}$ is nonzero.  Such a weight does not descend, because
$z_u-z_v\in\ker q$ is not homogeneous for it.  This remains a valid
counterexample in characteristic two, where the two distinct powers
still have nonzero coefficients.  Alternatively give both paths label
$1$ in characteristic zero: the original polynomial is $2x\ne0$,
but every weight induced from $x\mapsto t^{\omega}$ ties its two paths.
Thus compatible minimum-path isolation is not even necessary for
nonzero scalar output, while arbitrary path isolation is not sufficient
to certify nonzero output after the quotient. \hfill$\square$

\paragraph{What is refuted.}
The refuted assertion is the direct identification of a Hoza seed with
an ROABP variable-weight assignment.  Neither polynomial ROABP hitting
sets nor their possible derivation by a different construction is
refuted by this example.

\subsection{I-03. Three exact characterizations of basis isolation}
\label{I-03}

\paragraph{Statement.}
Let $S\subset\mathbb Z_{\geq0}^n$ index distinct monomials and let
$V=\operatorname{span}_F\{v_a:a\in S\}$ have dimension $r$ over an arbitrary field.
For $\omega\in\mathbb Z^n$ set $w_a=\langle\omega,a\rangle$.
For integer weights we use the same order-theoretic definition as
I-B02, allowing negative variable weights; the hitting-set
applications below use nonnegative weights.
The following are equivalent.
\begin{enumerate}
\item $\omega$ is a basis-isolating weight assignment in the sense of
this extended definition.
\item Every nonzero linear functional $\lambda\in V^*$ makes
$\sum_a\lambda(v_a)\mathbf x^a$ have a unique monomial of minimum weight.
\item At each weight level $c$, at most one coefficient $v_a$ with
$w_a=c$ survives modulo
$V_{<c}=\operatorname{span}\{v_b:w_b<c\}$.
\end{enumerate}
For a fixed coefficient basis $B$, write uniquely
$v_a=\sum_{b\in B}\lambda_{ab}v_b$.  It is the basis isolated by
$\omega$ exactly when
\[
 \langle\omega,a-b\rangle>0
 \quad(a\notin B,\ \lambda_{ab}\ne0),
 \qquad
 \langle\omega,b-b'\rangle\ne0\quad(b\ne b'\in B).
\tag{I.3}\label{eq:I-cone}
\]
The first inequalities alone characterize $B$ as the unique
minimum-total-weight coefficient basis.

\paragraph{Proof.}
Assume (1), with isolated basis $B$.  For a nonzero $\lambda$, take
the least-weight basis member $b$ on which $\lambda$ is nonzero.
It is unique.  Every nonbasis coefficient of no greater weight lies
in the span of smaller basis coefficients and is annihilated by
$\lambda$.  This proves (2), including its stronger unique-monomial
assertion.

If (3) fails, choose two nonzero coefficient classes at a level $c$.
There is a functional on $V/V_{<c}$ nonzero on both: for independent
classes assign value $1$ to both and extend a basis; for dependent
nonzero classes use any functional nonzero on one of them.  Pull it
back to $V$.  It vanishes below $c$ and is nonzero on two monomials
at $c$, contradicting (2).  This argument works also over finite
fields.  Thus (2) implies (3).

Under (3), choose the unique surviving monomial at every level where
the coefficient span grows.  These coefficients form a basis, one
per growth level.  All remaining coefficients at a level already
belong to the span of smaller levels.  Iterating through the finitely
many levels proves (1).  Unique coordinates in a fixed basis prove
the equivalence with (I.3).

If $B$ is the unique minimum-total-weight basis and
$\lambda_{ab}\ne0$, then $B-b+a$ is also a basis.  Its total weight
is strictly larger, so $w_a>w_b$.  Conversely, suppose these strict
inequalities hold.  Expand any other basis $C$ in the basis $B$.
Its change-of-basis determinant is nonzero, so at least one
permutation in its determinant expansion is nonzero.  Every common
member of $B\cap C$ is forced to match itself; every member
$a\in C\setminus B$ matches some $b\in B\setminus C$ with
$\lambda_{ab}\ne0$.  Summing the strict inequalities over this
matching gives $\sum_{a\in C}w_a>\sum_{b\in B}w_b$.
Thus $B$ is the unique minimum-total-weight basis. \hfill$\square$

\paragraph{Geometric consequence.}
For a fixed ordering of the weights within $B$, (I.3) is an open
rational polyhedral cone in the cocharacter space $\mathbb R^n$.
The union of these cones is precisely the set of BIWAs.  Equivalently,
the one-parameter subgroup simultaneously gives a monomial initial
form for every nonzero scalar projection of $D$.  Isolating one path
or one preselected scalar output has weaker quantifiers.

\subsection{I-04. Projected matroid vertices and polynomial random weights}
\label{I-04}

\paragraph{Statement.}
Let $\mathcal B$ be the coefficient bases of the data in I-03, and put
\[
 \beta(B)=\sum_{a\in B}a,\qquad
 Q=\operatorname{conv}\{\beta(B):B\in\mathcal B\}\subset\mathbb R^n.
\]
Every vertex of $Q$ is the image of exactly one indexed basis.
Assume $S\subseteq\{0,\ldots,\delta\}^{n}$, $r\ge1$, $\delta\ge1$.
Choose independently
$a_i\in\{1,\ldots,M_1\}$ and $b_i\in\{1,\ldots,M_2\}$.
With $L=n\delta M_2+1$, the weights
\[
 \omega_i=L a_i+b_i
\tag{I.4}
\]
are basis isolating with probability at least
\[
 1-\frac{nr\delta}{M_1}-\frac{\binom r2}{M_2}.
\tag{I.5}
\]
For example $M_1=4nr\delta$ and $M_2=4\max(1,r^2)$ give a positive
constant success probability and polynomially bounded weights.  This
is a randomized existence statement, not an explicit small family.

\paragraph{Proof of the vertex assertion.}
Suppose two distinct bases $B,C$ have the same image $\beta$.
Take $e\in B\setminus C$.  Let $\ell$ be the coordinate functional
of $v_e$ in the basis $B$, and expand $v_e=\sum_{f\in C}c_fv_f$.
Because $1=\ell(v_e)=\sum c_f\ell(v_f)$, some
$f\in C\setminus B$ has both $c_f\ne0$ and $\ell(v_f)\ne0$.
Hence $B-e+f$ and $C-f+e$ are both bases.  Their images are
$\beta+(f-e)$ and $\beta-(f-e)$, which are distinct because the
ground-set exponents are distinct.  Their midpoint is $\beta$,
so $\beta$ is not a vertex.  Every vertex has at least one preimage
from the finite defining set, proving the claim.

\paragraph{Proof of the probability bound.}
First consider any finite set $T\subseteq\{0,\ldots,D\}^{n}$ and
uniform independent weights $a_i\in[M_1]$.  Fix all weights except
$a_i$.  The least weight as a function of a real parameter $t=a_i$
is the lower envelope of at most $D+1$ affine functions, one for
each possible $i$-th coordinate.  Its slopes are among
$0,\ldots,D$, and there are at most $D$ breakpoints at which
different slopes minimize.  If two different points of $T$ minimize,
they differ in some coordinate $i$, and that coordinate weight is
one of these breakpoints.  Conditioning on the other coordinates
and applying a union bound gives probability at most $nD/M_1$.
Apply this with $T=\{\beta(B)\}$ and $D=r\delta$.
With probability at least $1-nr\delta/M_1$, there is a unique
minimizing point of $Q$.  It is a vertex and, by the preceding
argument, has a unique basis preimage $B$.

Conditional on the first weights and this $B$, two different
monomials in $B$ have equal $b$-weight with probability at most
$1/M_2$: condition on every $b_i$ but one whose exponent difference
is nonzero.  There is at most one possible value of that remaining
integer weight.  A union bound over pairs costs
$\binom r2/M_2$.  By I-03, the unique first-weight minimum basis
satisfies $a\cdot m>a\cdot b$ whenever the $b$-coordinate of
$v_m$ is nonzero and $m\notin B$.  Each difference is an integer
at least one.  The perturbation changes it by at most
$n\delta M_2$, so multiplication by $L=n\delta M_2+1$ preserves
strict positivity.  It also preserves inequalities between basis
members with unequal first weights, and separates any remaining
ties by the second weights.  Equation (I.3) proves (I.5).
All randomness and weights in this proof are integers; no assumption
on the characteristic of the coefficient field enters the argument.
\hfill$\square$

\paragraph{The remaining derandomization problem.}
The desired generator must hit these cones for coefficient matroids
arising from products of univariate matrices, with a family and
maximum weights both polynomial in $(n,w,\delta)$.  Hoza's theorem
hits path-polytope cones in independent vertex coordinates.  I-02
shows why those cocharacter spaces cannot be silently identified.
I-04 narrows the required object to normal cones of a specified
projected representable-matroid polytope, with the finite
within-basis tie hyperplanes removed.

\subsection{I-05. Positive Gram determinants certify bases, not compressed rank}
\label{I-05}

\paragraph{Statement.}
For complex coefficient vectors, choose coordinates $V=\mathbb C^r$
and form
\[
 G_D(\mathbf z)=\det\left(\sum_{a\in S}v_av_a^*\mathbf z^a\right).
\]
Then
\[
 G_D(\mathbf z)=\sum_{B\in\mathcal B}
                  |\det C_B|^2\mathbf z^{\beta(B)},
\qquad \operatorname{Newt}(G_D)=Q,
\tag{I.6}
\]
where $C$ has columns $v_a$.  At a vertex of $Q$, the coefficient
is the squared volume of its unique basis from I-04.  Nevertheless,
positivity of $G_D(t^\omega)$ does not imply that the coefficient
span of $D(t^\omega)$ retains its dimension.

\paragraph{Proof.}
Apply the Cauchy--Binet identity to
$C\operatorname{diag}(\mathbf z^a)C^*$.
Every nonzero summand is strictly positive in its coefficient, so
there are no cancellations between bases with the same exponent
sum; their convex hull is exactly the Newton polytope.  I-04
identifies its vertex coefficients.  For the negative assertion,
take $D(x,y)=e_1x+e_2y$ and $\omega=(1,1)$.  Its Gram determinant
is $xy$ and specializes to $t^2$, whereas
$D(t,t)=t(e_1+e_2)$ has coefficient span of dimension one instead
of two.  The unique coefficient basis has two equal member weights;
this is exactly the extra condition in I-03. \hfill$\square$

\paragraph{Field and geometric qualification.}
The conjugate transpose and positivity in (I.6) are statements on
a complex vector space with a Hermitian structure.  They are not
positivity statements in an arbitrary characteristic.  The
determinantal/Cauchy--Binet identity itself has an integral
two-matrix form with independent dual variables.  Specializing the
second matrix to a conjugate transpose is a real, not a complex
algebraic, operation.  This qualification is essential in translating
the formula into a scheme over $\mathbb Z$.

\subsection{I-06. Exact black-box implementation of nonlinear exponent weights can be exponential}
\label{I-06}

\paragraph{Statement.}
Let $K$ be any field, let $t$ be transcendental over it, and define
a linear functional on polynomials of individual degree at most $2d$
by
\[
 \mathcal L(\mathbf x^\alpha)=\prod_{i=1}^n h_i(\alpha_i),
 \qquad h_i(e)\in K(t).
\]
Put $H_i=(h_i(a+b))_{0\le a,b\le d}$.
If one exact universal implementation is
\[
 \mathcal L(f)=\sum_{j=1}^M c_j f(\xi_j),
\tag{I.7}
\]
where the points and coefficients are fixed independently of $f$
over any extension of $K(t)$, then
\[
 M\ \ge\ \prod_{i=1}^n\operatorname{rank}(H_i).
\tag{I.8}
\]
In particular $\mathcal L(\mathbf x^\alpha)=
t^{\sum_i\alpha_i^2}$ requires $M\ge(d+1)^n$.
For degree two, the nonlinear exponent filters supplied by every
Hoza seed that separates three exponent symbols also have the
lower bound $M\ge2^n$.

\paragraph{Proof.}
Index a moment matrix by
$\alpha,\beta\in\{0,\ldots,d\}^n$ and set
$H_{\alpha\beta}=\mathcal L(\mathbf x^{\alpha+\beta})$.
It is the tensor product $\bigotimes_iH_i$, with rank the product
of the ranks.  Each evaluation in (I.7) contributes the rank-one
matrix $c_j(\xi_j^\alpha)_\alpha(\xi_j^\beta)_\beta^{\mathsf T}$,
so its rank is at most $M$.  Extension of scalars does not change
the rank, proving (I.8).

For $h_i(e)=t^{e^2}$, factor the $(a,b)$ entry as
$t^{a^2}t^{b^2}(t^{2a})^b$.  The remaining matrix is Vandermonde
on the distinct elements $1,t^2,\ldots,t^{2d}$ and has full rank,
in every characteristic.

For the Hoza assertion take three distinct binary codewords
$C(0),C(1),C(2)$ and put
$h_x(e)=\sum_j6^{\ell-j}C(e)_j x_j$ as integers.
If $h_x(0)+h_x(2)=2h_x(1)$, then
\[
 \sum_j6^{\ell-j}x_j\bigl(C(0)_j+C(2)_j-2C(1)_j\bigr)=0.
\]
The digits lie between $-2$ and $2$.  The first nonzero digit
dominates the absolute sum of all later digits, since
$2\sum_{k=0}^{q-1}6^k<6^q$.  Thus every digit vanishes, which for
binary digits says the three selected code bits agree in every
selected coordinate.  Hence all three values $h_x(e)$ coincide.
Conversely constant values give that equality.  Any seed isolating
the three-arm graph is nonconstant on its three internal vertices,
so the equality fails.  At layer $i$, the scalar multiplier
$6^{\nu_2(i)}$ in (I.1) preserves this failure.  Consequently the
$2\times2$ Hankel determinant of the exponent filter
\[
 e\longmapsto t^{6^{\nu_2(i)}h_x(e)}
\]
is the difference of two distinct monomials, and has nonzero
determinant over $K(t)$.  Equation (I.8) now gives $2^n$.
\hfill$\square$

\paragraph{Scope of the lower bound.}
This rules out the proposed exact linear interpolation implementation
of independently reweighted monomial letters by a polynomial number
of evaluations of the original black box.  It is not a lower bound
on general ROABP PIT, adaptive reconstruction, or approximate
integration.  Even requiring (I.7) only for width-one ROABPs implies
the same identities on every monomial and hence the same lower
bound by linearity; width-one PIT itself remains elementary.

\subsection{I-07. Haar averaging repairs invariance by discarding state isolation}
\label{I-07}

\paragraph{Statement.}
For a fixed complex $m\times m$ matrix $D$, the function
$A\mapsto\operatorname{tr}(DA)$ is invariant under simultaneous
conjugation if and only if $D$ is scalar.  Moreover
\[
 \int_{U(m)}UDU^*\,dU=\frac{\operatorname{tr}D}{m}I.
\tag{I.9}
\]
Thus independently averaging fixed diagonal state-weight insertions
in a trace word replaces each insertion by a scalar and recovers
a scalar multiple of the original trace word.  If the same unitary
is used for two insertions, the exact formula for $m\ge2$ is
\[
 \int\operatorname{tr}(A UDU^* B UEU^*)\,dU
 =\alpha\operatorname{tr}(AB)+\beta\operatorname{tr}A\operatorname{tr}B,
\tag{I.10}
\]
where
\[
 \alpha=\frac{m\operatorname{tr}D\operatorname{tr}E-
                    \operatorname{tr}(DE)}{m(m^2-1)},\qquad
 \beta=\frac{m\operatorname{tr}(DE)-
                    \operatorname{tr}D\operatorname{tr}E}{m(m^2-1)}.
\]
Even positive semidefinite $D,E$ need not make these coefficients
nonnegative: for orthogonal rank-one projections,
$\beta=-1/[m(m^2-1)]$.

\paragraph{Proof.}
Conjugation invariance for every $A$ gives
$\operatorname{tr}((g^{-1}Dg-D)A)=0$ for all $A$.  The trace pairing
is nondegenerate, so $D$ commutes with every invertible matrix and
is scalar.  The reverse direction is immediate.  Haar invariance
makes the left side of (I.9) commute with every unitary; the
commutant is scalar, and its trace determines the scalar.  Repeated
integration gives the independent-insertion claim.

For completeness, in (I.10) define
$T=\int(UDU^*)\otimes(UEU^*)\,dU$.  The commutant of
$U\otimes U$ is spanned by the identity $I$ and tensor flip $S$.
Here is an elementary verification: commutation with diagonal unitaries
forces matrix entries to preserve the unordered pair of tensor indices;
permutation unitaries make the two-dimensional off-diagonal blocks
$\left(\begin{smallmatrix}b&c\\c&b\end{smallmatrix}\right)$ and
the diagonal blocks a common scalar $a$.  A rotation mixing two basis
vectors forces $a=b+c$.  The resulting operator is $bI+cS$.
Thus $T=\alpha I+\beta S$.  Its trace and its trace after multiplication
by $S$ are respectively
\[
 m^2\alpha+m\beta=\operatorname{tr}D\operatorname{tr}E,
 \qquad m\alpha+m^2\beta=\operatorname{tr}(DE).
\]
Solving gives the stated values.  Finally
$\operatorname{tr}(A X B Y)=
\operatorname{tr}((A\otimes B)(X\otimes Y)S)$;
contracting $T$ gives (I.10).  For orthogonal rank-one projections,
$\operatorname{tr}D=\operatorname{tr}E=1$ and
$\operatorname{tr}(DE)=0$, proving the sign assertion.
\hfill$\square$

\paragraph{Meaning for the proposed transfer.}
Hoza's internal-state weights enter matrix products as fixed
diagonal insertions.  Their lack of conjugation invariance cannot
be repaired by independent Haar averaging while preserving the
state information.  Correlated averaging changes the trace word
into a combination of different trace invariants and can reintroduce
negative coefficients.  Neither averaging is the missing
basis-isolation construction.

\subsection{I-08. The exact matrix-invariant consequence, conditional on the remaining gate}
\label{I-08}

\paragraph{Statement.}
Suppose one has an explicit hitting set $\mathcal H\subset F^{m^2}$
of polynomial size, constructible in polynomial arithmetic time,
for ROABPs of depth $m^2$, width at most $2m^2$, and individual
degree less than $r$.  Let $F$ have characteristic zero and put
\[
 M(x)=\sum_{j=0}^{r-1}M_jx^j,\qquad
 T_{\ell,\alpha}(\mathbf M)=
 \operatorname{tr}\bigl(M(\alpha_1)\cdots M(\alpha_\ell)\bigr),
 \quad1\le\ell\le m^2,quad\alpha\in\mathcal H.
\tag{I.11}
\]
These at most $m^2|\mathcal H|$ invariants are an explicit
polynomial-size separating family for simultaneous conjugation.
Over algebraically closed $F$, the invariant ring is finite over
the subring they generate.  This is precisely the specialization
of Forbes and Shpilka (2013), Construction~3.7, Theorem~3.8 and
Corollary~3.10, to a polynomial hitting set.  It supplies the
matrix-invariant instance of explicit Noether normalization in
the sense of that paper.  It does not by itself extract an
algebraically independent system of parameters of cardinality
equal to the quotient dimension.

\paragraph{Derivation and exact source dependence.}
The coefficient of $x_1^{a_1}\cdots x_\ell^{a_\ell}$ in
$\operatorname{tr}(M(x_1)\cdots M(x_\ell))$ is the trace word
$\operatorname{tr}(M_{a_1}\cdots M_{a_\ell})$.  The latter words,
for $\ell\le m^2$, generate the characteristic-zero invariant ring
as recorded in Forbes--Shpilka Theorem~1.3 and Corollary~3.6.
For two tuples $A,B$, the difference of their trace-generating
polynomials is an ROABP of width at most $2m^2$, by their Lemma~3.2.
The hitting-set assumption therefore detects exactly whether
some trace generator separates $A$ and $B$.  Their Theorem~3.9
(the stated specialization of Derksen--Kemper Theorem~2.3.12)
then gives integrality, and finite generation upgrades this to
a finite morphism.  The only new hypothesis in this specialization
is the polynomial hitting-set gate; it has not been proved here.
\hfill$\square$

\paragraph{Consequences that do not follow.}
This does not prove general arithmetic PIT, arbitrary explicit
Noether normalization, permanent circuit lower bounds, or
$P\ne NP$.  Forbes--Shpilka already prove deterministic NC
orbit-closure intersection in Theorem~1.15 using white-box PIT.
Their discussion after Theorems~1.13--1.15 explicitly separates
the matrix-invariant case from the stronger general explicit-variety
problem.  In the supplied revised GCT~V source, the general-PIT
equivalence is Theorem~5.14, and the matrix-specific ROABP replacement
is \S7.5, Lemma~7.9.  The earlier GCT~V theorem numbering cited by
Forbes--Shpilka must not be transferred to the revised source.  No new polynomial hitting set
may be inserted into (I.11) until its own gate is closed.

\subsection{I-09. An exact positive transfer-matrix certificate remains white box}
\label{I-09}

\paragraph{Statement.}
For a complex ROABP in known matrix-product form
\[
 f(\mathbf x)=u^{\mathsf T}\prod_{i=1}^n
                  \left(\sum_{a=0}^{\delta}A_{i,a}x_i^a\right)v
             =\sum_{\alpha}c_{\alpha}\mathbf x^{\alpha},
\]
its coefficient squared-norm generating polynomial is
\[
 \sum_{\alpha}|c_\alpha|^2\mathbf z^\alpha
 = (u\otimes\bar u)^{\mathsf T}
     \prod_{i=1}^n\left(\sum_{a=0}^{\delta}
                    A_{i,a}\otimes\overline{A_{i,a}}\,z_i^a\right)
                     (v\otimes\bar v).
\tag{I.12}
\]
In particular $f=0$ if and only if the right side at
$\mathbf z=\mathbf1$ is zero.  Its matrices have width $w^2$.
The equality also gives the usual torus identity
\[
 \int_{[0,2\pi]^n}|f(e^{i\theta_1},\ldots,e^{i\theta_n})|^2
                \prod_i\frac{d\theta_i}{2\pi}
       =\sum_\alpha|c_\alpha|^2.
\tag{I.13}
\]

\paragraph{Proof.}
The coefficient $c_\alpha$ is
$u^{\mathsf T}A_{1,\alpha_1}\cdots A_{n,\alpha_n}v$.
Multiplicativity of tensor products identifies its squared
absolute value with the product of the corresponding tensor
matrices and boundary vectors.  Summing independently over each
$\alpha_i$ gives (I.12).  The sum is a finite sum of nonnegative
real numbers, and it is zero exactly when every coefficient is
zero.  In (I.13), expansion of the square and integration of
$e^{i\langle\alpha-\beta,\theta\rangle}$ kills exactly the terms
$\alpha\ne\beta$. \hfill$\square$

\paragraph{Why this does not close the black-box gate.}
The transfer matrices use the individual coefficient matrices
$A_{i,a}$, which a black box does not supply.  The torus integral
is another exact expression, but the product quadrature that is
exact for all such degree bounds uses $(\delta+1)^n$ nodes.
No lower bound on all specialized ROABP quadratures is asserted.
I-06 proves a lower bound for the different, explicitly specified
nonlinear exponent filter.  Positivity in (I.12) closes the
white-box zero-testing question only; it does not supply the
missing uniform polynomial set of original evaluations.

\subsection{I-10. Schemes, cocharacters, flags, and flat limits}
\label{I-10}

\paragraph{Statement and dictionary.}
Let $R$ be a domain carrying the relevant coefficients, and let
$T=\mathbb G_{m,R}^{n}$.  The integer weight vector $\omega$
is the cocharacter
$\lambda_\omega:\mathbb G_{m,R}\to T$,
$t\mapsto(t^{\omega_1},\ldots,t^{\omega_n})$.
For nonzero $f=\sum_a c_a\mathbf x^a$ and
$\mu=\min_{c_a\ne0}\langle\omega,a\rangle$, define
\[
 \widetilde f(t,\mathbf x)=
 \sum_a c_a t^{\langle\omega,a\rangle-\mu}\mathbf x^a.
\tag{I.14}
\]
This is a polynomial over $R[t]$, its special fiber is the
weight-initial form, and its generic fiber is the original
polynomial under a torus change of coordinates.  A unique
minimizing exponent gives a monomial special fiber.  Over a field,
the hypersurface family defined by (I.14) is flat over the
$t$-line.  For a homogeneous ideal, the corresponding weight
homogenization is the usual flat Rees/initial-ideal degeneration;
this is stronger than merely selecting a least monomial of one
polynomial and must not be conflated with it.

For a vector polynomial, I-03 identifies a complete coefficient
flag whose successive quotients have one supported weight
monomial.  With an explicit coefficient matrix, nonvanishing
maximal minors define a matroid stratum of the appropriate
Grassmannian.  The exponents of its Pl\"ucker coordinates under
$\lambda_\omega$ are the sums $\beta(B)$.  Its torus orbit has
weight polytope $Q$; I-04 says a vertex with nonzero Pl\"ucker
coordinate corresponds to a unique indexed coefficient basis.
The BIWA cones refine its normal fan by the individual
within-basis weight comparisons in (I.3).

\paragraph{Proof of the flatness assertion and the dictionary.}
Formula (I.14) has no negative power of $t$ and is not divisible
by $t$.  Over a field $F$, for every irreducible polynomial
$p(t)$ the image of $\widetilde f$ modulo $p(t)$ is nonzero:
for $p=t$ use its nonzero initial form; otherwise $t$ is invertible
in $F[t]/(p)$ and its monomial coefficients remain nonzero.
Thus no nonzero $p(t)$ divides $\widetilde f$ in the UFD
$F[t,\mathbf x]$.  If $p(t)g\in(\widetilde f)$, primality of
the irreducible factors of $\widetilde f$ implies
$g\in(\widetilde f)$.  Hence the quotient by
$(\widetilde f)$ is torsion-free over the PID $F[t]$, and
therefore flat.  The generic-fiber statement follows by
inverting $t$ in (I.14).  Exterior products of coefficient
columns are precisely Pl\"ucker coordinates, and their weights
add, giving $\beta(B)$.  Equation (I.3) then identifies the
normal-cone refinement.  This proves each part of the stated
dictionary directly. \hfill$\square$

\paragraph{Integral qualification.}
The torus and monomial maps are defined over $\mathbb Z$.
A specified coefficient matroid and its rank flag persist on an
open localization where the finitely many designated nonzero
minors remain nonzero; they need not persist in every residue
characteristic.  Over a general base $R$, flatness of a
hypersurface quotient also needs the relevant coefficient and
fiber hypotheses.  The field proof above is not a claim of
automatic flatness over arbitrary $R$.  The Grothendieck locator
table in the main paper should be used for the exact SGA/FGA/EGA
placement of these notions, rather than inventing a theorem
number for these elementary derived statements.

\subsection{I-11. Exterior and subspace-state repairs do not yet produce a graph generator}
\label{I-11}

\paragraph{Checked finite assertions.}
For coefficient ambient dimension $k$, the standard coordinate
space $\bigwedge^rF^k$ has dimension $\binom kr$.  Over
$\mathbb F_q$, there are at least $q^{r(k-r)}$ distinct
$r$-dimensional subspaces of $F_q^k$: the graphs of the
$q^{r(k-r)}$ linear maps $F_q^r\to F_q^{k-r}$ are all distinct.
Thus a literal graph with one state for each such subspace is
not polynomial in $k$ when $r$ grows with $k$.

Exterior coordinates also retain cancellation.  For two equal
vectors $v_1=v_2=e_1+e_2$, the determinant is
$1\cdot1-1\cdot1=0$, although its expansion has two nonzero
permutation terms.  Making one term uniquely minimum by
path-dependent weights changes the determinant rather than
certifying its nonvanishing.  Tensoring $r$ independently chosen
word trajectories gives a standard realization with $w^{O(r)}$
states; for fixed $r$ this removes that particular size cost,
but the cancellation and weight-descent problems remain.

\paragraph{Proof and scope.}
The exterior dimension follows from the basis wedges indexed
by $r$-element subsets of a fixed $k$-element basis.  The graph
of a linear map is an $r$-dimensional subspace, and projection
onto the first summand recovers the map, proving distinctness.
The determinant example is immediate.  The tensor product
construction records one of $w$ state positions for each of
the $r$ trajectories, so has $w^r$ positions before additional
boundary coordinates; determinant signs are imposed at the
boundary and do not disappear.  These are size and correctness
facts about the specified constructions, not lower bounds on
all possible reductions.  A scalar output can itself have
coefficient-space dimension one, so constant coefficient rank
alone cannot make the general ROABP hitting-set task disappear.
\hfill$\square$

\subsection{I-12. Resolving coefficient support is an arithmetic counting gate}
\label{I-12}

\paragraph{Background I-B05.}
Allender and Ogihara (1996), \S2, p.~4, define the accepting-minus-rejecting
gap model; \S5, Definition~2, p.~11, defines $C_{=}L$ by the vanishing
of a $\operatorname{GapL}$ function.  Their Theorem~1, p.~2, identifies
$\operatorname{GapL}$ with logarithmic-space reductions to integer
determinant, and Theorem~14, p.~12, gives the corresponding singularity
completeness result.  The following ROABP formulation is proved directly.

\paragraph{Statement.}
Restrict to explicitly given layered ROABPs over $\mathbb Z$
(equivalently over $\mathbb Q$) with coefficients in
$\{0,1,-1\}$ and polynomially bounded width and length.  The input
lists the coefficient entries and the queried exponent vector;
univariate edge labels may be given densely or as canonical sparse
lists with binary exponents.  Deciding
whether a specified monomial coefficient is nonzero is complete,
under deterministic logarithmic-space reductions, for the class
\[
 \operatorname{coC}_{=}L
   =\{\{x:g(x)\ne0\}:g\in\operatorname{GapL}\}.
\]
Here $\operatorname{GapL}$ is defined by the difference of the
numbers of accepting and rejecting paths of polynomial-time
nondeterministic logarithmic-space machines.  Hardness already
holds for ROABPs whose polynomial is one monomial times a scalar.
Therefore a logarithmic-space uniform polynomial-size reduction
which represents exactly nonzero coefficient support by ordinary
directed reachability would imply
$\operatorname{coC}_{=}L\subseteq NL$.

\paragraph{Proof.}
For membership, guess a state path in the layered ROABP.  At layer
$i$ choose the term of the prescribed exponent from its univariate
edge label.  A zero term contributes nothing; otherwise track the
parity of the negative signs.  The difference between positive and
negative paths is exactly the specified monomial coefficient.
State indices, a layer counter, and a sign bit require logarithmic
space.  Neutral or invalid branches can be arranged to contribute
one accepting and one rejecting path, hence zero gap.

For hardness, start with a polynomial-time nondeterministic
logarithmic-space machine computing a gap $g(x)$.  Its time-unrolled
configuration graph has polynomial width and length and is
constructible from $x$ in logarithmic space.  Pad halted computations
by deterministic transitions, retaining their accepting or rejecting
tag, so every path has the same length.  Mark the final edge
of accepting paths by $1$ and of rejecting paths by $-1$;
all preceding edges have label $1$.  Separate parallel transitions
by an intermediate layer if necessary.  Multiply every edge label
in layer $i$ by a fresh variable $x_i$.  Every source-to-target path
then has the same variable monomial $\prod_i x_i$, with its
accepting or rejecting sign, so the output is
$g(x)\prod_i x_i$.  This is a known-order multilinear ROABP with
coefficients $0,1,-1$.  Its unique possible monomial coefficient
is nonzero exactly when $g(x)\ne0$.  Finally, composition with the
postulated graph reduction puts every such language in $NL$.
\hfill$\square$

\paragraph{Relation to the other gates.}
This is a precise complexity formulation of the operation which
would have to cancel equal-monomial paths before graph isolation.
It is not a separation of complexity classes.  Nor does it rule
out a different polynomial-time algebraic construction: the
explicit ROABP supplies its matrices, and white-box polynomial
identity testing is already possible.  In scheme language the
coefficient is a regular function of the edge coefficients;
testing its vanishing uses the kernel of a coefficient map, not
merely whether its parameterizing path set is empty.  The diamond
in I-02 is the smallest instance of this distinction.

\subsection{I-13. Exact universal energy cubature needs the full coefficient dimension}
\label{I-13}

\paragraph{Statement.}
Let $n\ge1$, $\delta\ge0$, and $D=(\delta+1)^n$.  Fix points
$\xi_1,\ldots,\xi_Q\in\mathbb C^n$ and a complex $Q\times Q$ matrix $C$,
independent of the polynomial.  Suppose that every width-one product
\[
 f(\mathbf x)=\prod_{i=1}^n\left(\sum_{a=0}^{\delta}z_{i,a}x_i^a\right),
 \qquad z_{i,a}\in\mathbb C,
\]
satisfies the exact identity
\[
 \sum_{\alpha\in\{0,\ldots,\delta\}^n}|[\mathbf x^\alpha]f|^2
  =\sum_{j,k=1}^Q C_{jk}\,\overline{f(\xi_j)}f(\xi_k).
\tag{I.13}\label{eq:I-cubature}
\]
Then $Q\ge D$.  In particular, signed cubature
$\sum_j c_j|f(\xi_j)|^2$ requires at least $D$ nodes.  More generally,
a formula with $M$ fixed paired terms
$\sum_{j=1}^M c_j\overline{f(\xi_j)}f(\eta_j)$ requires $M\ge D$.
The bound is attained for the diagonal formula by the Cartesian product
of the $(\delta+1)$st roots of unity in each coordinate, all with weight
$1/D$.

\paragraph{Proof.}
Index coordinates by $\alpha\in\{0,\ldots,\delta\}^n$ and put
$a_\alpha=\prod_i z_{i,\alpha_i}$.  Let the evaluation matrix be
$E_{j\alpha}=\xi_j^\alpha$.  The asserted identity is
\[
 a^*(I_D-E^*CE)a=0
\quad\text{for every decomposable coefficient vector }
a=\bigotimes_i(z_{i,0},\ldots,z_{i,\delta}).
\]
Expanding gives a polynomial identity in the variables
$z_{i,a}$ and their conjugates.  Its monomials
\[
 \prod_i\overline{z_{i,\alpha_i}}z_{i,\beta_i},
 \qquad (\alpha,\beta)\in\{0,\ldots,\delta\}^{2n},
\]
are linearly independent as functions of the complex variables.
Indeed, writing each $z=u+iv$ identifies the polynomial ring in
formal $z,\bar z$ with the polynomial ring in $u,v$ by an invertible
complex linear change of variables; a polynomial vanishing on every
real tuple $(u,v)$ has every coefficient zero.  Distinct pairs
$(\alpha,\beta)$ give distinct displayed monomials.  Consequently
$I_D=E^*CE$, so
$D=\operatorname{rank}(I_D)\le\operatorname{rank}(E)\le Q$.
For paired terms the same coefficient comparison represents $I_D$
as a sum of $M$ rank-at-most-one matrices, proving $M\ge D$.

Finally, averaging $\xi^\alpha\overline{\xi^\beta}$ over the stated
root-of-unity grid gives $1$ when $\alpha=\beta$ and $0$ otherwise:
in each coordinate the difference lies between $-\delta$ and
$\delta$, whose only multiple of $\delta+1$ is zero.  Expanding
$|f|^2$ therefore proves the matching formula, for every polynomial
of the stated individual degree, including the width-one class.
\hfill$\square$

\paragraph{Scope and consequence for I-09.}
Torus orthogonality identifies the left side of (I.13) with
$\int_{(S^1)^n}|f(z)|^2\,dz$.  Thus the positive white-box energy
computation in I-09 cannot be replaced by a polynomial number of
fixed black-box samples through any exact quadratic formula of the
displayed kind, even for width one and fixed $\delta\ge1$ as $n$ grows.
This is a lower bound for universal exact cubature, not for hitting
sets: for width one the substitution $x_i=t$ gives the nonzero
univariate product $\prod_i f_i(t)$ of degree at most $n\delta$.
Evaluation at $n\delta+1$ distinct scalars already gives a polynomial
hitting set for that class.  The rank argument makes no claim about
adaptive procedures or other nonlinear uses of evaluations.

\subsection{I-14. Joint compression is a Grassmannian transversality condition}
\label{I-14}

\paragraph{Statement.}
Write $D=\sum_{a\in S}v_a\mathbf x^a$, choose coordinates on
$V=\operatorname{span}\{v_a\}$, and let $C=(v_a)_{a\in S}$ be the
resulting full-row-rank $r\times s$ matrix, where $s=|S|$.
For a finite family $\mathcal W$ of integer weights, form a matrix
$P_{\mathcal W}$ with rows indexed by $S$ and columns indexed by
the pairs $(\omega,c)$ that occur as a monomial weight.  Substitutions
are interpreted as Laurent polynomials if necessary.  Set
\[
 (P_{\mathcal W})_{a,(\omega,c)}
   =\begin{cases}1&\langle\omega,a\rangle=c,\\0&\text{otherwise.}\end{cases}
\]
The following conditions are equivalent:
\begin{enumerate}
\item the coefficients of all the substitutions
$D(t^{\omega_1},\ldots,t^{\omega_n})$, taken jointly, span $V$;
\item $\operatorname{rank}(CP_{\mathcal W})=r$;
\item $\operatorname{im}(C^{\mathsf T})\cap
\ker(P_{\mathcal W}^{\mathsf T})=\{0\}$;
\item for every nonzero $\lambda\in V^*$, at least one substituted
scalar polynomial $\lambda(D(t^\omega))$ is nonzero.
\end{enumerate}
Equivalently, for some $r$-element set $T$ of columns of $P_{\mathcal W}$,
\[
 \sum_{\substack{B\subseteq S\\|B|=r}}
   \det(C_B)\det((P_{\mathcal W})_{B,T})\ne0.
\tag{I.14}\label{eq:I-joint-Plucker}
\]
Thus the failure locus on $\operatorname{Gr}(r,F^s)$ is the
intersection with the linear equations obtained by setting every
expression in (I.14) to zero.  These depend on Pl\"ucker coordinates,
including their values, rather than only on the coefficient matroid.

\paragraph{Proof.}
The $(\omega,c)$ column of $CP_{\mathcal W}$ is exactly
$\sum_{\langle\omega,a\rangle=c}v_a$, the coefficient at that
weight after substitution.  This proves (1)$\Leftrightarrow$(2).
Since $C^{\mathsf T}:F^r\to F^s$ is injective, the transpose of
$CP_{\mathcal W}$ is injective exactly when the intersection in
(3) is zero.  A proper subspace of $V$ has a nonzero annihilating
linear functional, proving (1)$\Leftrightarrow$(4).  Finally,
Cauchy--Binet applied to each $r$-column submatrix of
$CP_{\mathcal W}$ gives (I.14).  These are linear forms in the
Pl\"ucker coordinates $\det(C_B)$, proving the geometric description.

Joint preservation can hold even though no individual member is
basis isolating.  Take
\[
 D=e_1(x_1-x_2)+e_2(x_2-x_3),\qquad
 \mathcal W=\{(0,0,1),(1,0,0)\}.
\]
The two substituted vector polynomials are respectively
$(1-t)e_2$ and $(t-1)e_1$.  Each has coefficient rank one and so
fails I-03, but their coefficients jointly span the original
two-dimensional space.  To see why the matroid alone is insufficient,
over a field of characteristic different from two compare coefficient
rows $(1,1)$ and $(1,-1)$ on the support $\{x_1,x_2\}$ under weight
$(1,1)$.  They have the same rank-one coefficient matroid, while
their compressed rows are respectively $2$ and $0$.
\hfill$\square$

\paragraph{A sufficient black-box gate.}
If $S\subseteq\{0,\ldots,\delta\}^n$, all weights are nonnegative
and at most $B$, and joint preservation holds, evaluating the
$|\mathcal W|$ curves at $n\delta B+1$ distinct field elements
hits every nonzero scalar projection of $D$.  Indeed, condition
(4) supplies a nonzero univariate polynomial of degree at most
$n\delta B$.  This exact sufficient condition is weaker than
requiring one BIWA.  It does not remove cancellation from the
linear Pl\"ucker expressions in (I.14).

\subsection{I-15. Bounded affine decodings of the seed masks have width-two obstructions}
\label{I-15}

\paragraph{Statement.}
Suppose a proposed family of variable weights has the form
\[
 \omega(s)=b+Ts,\qquad
 b\in\mathbb Z^n,\quad T\in\mathbb Z^{n\times k},
\]
where every entry of $b$ and $T$ has absolute value at most $M$.
The seed set can be any subset of $\mathbb Z^k$.  If
\[
 2^n>(2nM+1)^{k+1},
\tag{I.15}\label{eq:I-affine-mask-count}
\]
there are distinct $\alpha,\beta\in\{0,1\}^n$ such that the
nonzero multilinear binomial $\mathbf x^\alpha-\mathbf x^\beta$
vanishes identically on every curve $\mathbf x=t^{\omega(s)}$.
This binomial has a width-two ROABP in every variable order.
In particular the whole family fails even the joint condition I-14.

\paragraph{Proof.}
Map the $2^n$ Boolean exponent vectors $\alpha$ to
\[
 \left(b^{\mathsf T}\alpha,T^{\mathsf T}\alpha\right)\in
 \{-nM,\ldots,nM\}^{k+1}.
\]
By (I.15), two distinct vectors have the same image.  Then
$\langle\omega(s),\alpha-\beta\rangle=0$ for every seed $s$,
so the two substituted powers of $t$ coincide.  This is a Laurent
identity if negative weights are allowed, and a polynomial identity
when the weights are nonnegative.  Distinct monomials have nonzero
coefficients in their difference over any field, including
characteristic two.  The diagonal matrix product with factors
$\operatorname{diag}(x_i^{\alpha_i},x_i^{\beta_i})$ and boundaries
$(1,-1)$ and $(1,1)^{\mathsf T}$ computes the binomial.  Reordering
the commuting diagonal factors gives every variable order.
\hfill$\square$

\paragraph{A necessary dimension for any bounded weight family.}
More generally, let $\mathcal W\subseteq\mathbb Z^n$ have all
coordinates of absolute value at most $B\ge1$, and let
$k=\dim_{\mathbb Q}\operatorname{span}_{\mathbb Q}\mathcal W$.
If $\mathcal W$ hits every multilinear binomial by monomial curves,
then necessarily
\[
 k\ge\frac{n}{\log_2(2nB+1)}.
\]
Indeed, select $k$ members of $\mathcal W$ forming a rational basis
and apply the same pigeonhole argument to their integer dot products
with Boolean exponent vectors.  Equality on this basis implies
equality on every member of $\mathcal W$, independently of how
the seed was decoded.  Thus polynomially bounded output weights
computed as polynomials of any fixed degree $h$ in
$\ell=O(\log n)$ Boolean seed bits also fail for sufficiently large
$n$: their span has dimension at most
$\sum_{j=0}^h\binom{\ell}{j}=O((\log n)^h)$, below the necessary
bound.  The dimension estimate follows by grouping the coefficient
vectors of the multilinear seed monomials; on Boolean seeds powers
of a bit reduce to that bit.

\paragraph{Application to the actual mask formula.}
Every weight in Hoza's formula (I.1) is an integer linear form in
the common $\ell$ seed bits.  Hence assigning those forms to
variables, or taking any fixed linear averaging or decoding with
polynomially bounded integer coefficients, gives $\omega(s)=b+Ts$
with $k=\ell$ and polynomially bounded entries.  A rational average
with a common polynomially bounded denominator has the same property
after clearing that denominator.  Since $\ell=O(\log(WN))$, if
$W,N$ and the resulting entry bound are polynomial in the number
$n$ of variables, condition (I.15) holds for all sufficiently large
$n$.  This rules out this specified seed-to-variable decoding,
even if all its members are used jointly.

For the most direct alphabet decoding the counterexample is explicit
without counting.  Let
$H_s(a)=\sum_j6^{\ell-j}C(a)_js_j$.  Matching letters $0,1$ by
an affine exponent map at layer $i$ gives
\[
 \omega_i(s)=6^{\nu_2(i)}\bigl(H_s(1)-H_s(0)\bigr).
\]
Layers $1$ and $3$ receive the same weight, so $x_1-x_3$ vanishes
for every seed.  This uses the actual layer dependence in (I.1)
and holds for every choice of its code.  There is no contradiction
with minimum-path isolation: the graph weights themselves need not
descend to these variable weights.  The argument does not rule out
nonlinear decoders whose outputs span sufficiently many independent
directions, or unbounded integer weights.

\subsection{I-16. The finite-field row-state lift and its uniform evaluation limit}
\label{I-16}

\paragraph{Statement.}
Over $F=\mathbb F_q$, an explicitly given width-$w$ ROABP
\[
 f=u^{\mathsf T}\prod_{i=1}^n
       \left(\sum_{a=0}^{\delta}A_{i,a}x_i^a\right)v
\]
has an exact deterministic coefficient-support automaton with
at most $q^w$ states per layer.  A state is a row vector
$r\in F^{1\times w}$; its letter-$a$ transition at layer $i$ is
$r\mapsto rA_{i,a}$.  The initial state is $u^{\mathsf T}$, and
a final state is accepting exactly when $rv\ne0$.
Nevertheless, if $0\le\delta<q$, every base-field hitting set
$H\subseteq F^n$ for all width-one individual-degree-$\delta$
polynomials satisfies
\[
 |H|\ \ge\ \left(\frac{q}{q-\delta}\right)^n.
\tag{I.16}\label{eq:I-small-field-lower}
\]
For $\delta=q-1$ all $q^n$ points are necessary and sufficient.
Thus polynomial state-space size alone cannot produce a polynomial
uniform base-field hitting set, even when $w=1$ and $q$ is fixed.

\paragraph{Proof.}
On the word $(a_1,\ldots,a_n)$, the automaton's final scalar is
$u^{\mathsf T}A_{1,a_1}\cdots A_{n,a_n}v$, precisely the monomial
coefficient.  Matrix multiplication has already performed all path
cancellation before the state is tested.  To obtain a simple layered
graph, subdivide each transition with an intermediate vertex labelled
by its current row and its letter; append a common sink after all
accepting states.  The width is at most $(\delta+1)q^w$ and the
length is $2n+1$.  It is constructible from the coefficient matrices
in time polynomial in this displayed size and the field-operation
cost.  Hoza applies to this coefficient-dependent graph, with the
corresponding polynomial-in-$q^w$ seed family; its path weights
still depend on the row states.

For the lower bound, choose independently a $\delta$-element subset
$R_i\subseteq F$ for each coordinate and consider
\[
 f_{R_1,\ldots,R_n}(\mathbf x)
     =\prod_{i=1}^n\prod_{a\in R_i}(x_i-a).
\]
These are $\binom{q}{\delta}^{n}$ distinct nonzero width-one polynomials.
A fixed point $\xi\in F^n$ hits exactly
$\binom{q-1}{\delta}^n$ of them, since each $R_i$ must avoid
$\xi_i$.  Covering all these polynomials therefore requires
\[
 |H|\binom{q-1}{\delta}^n\ge\binom{q}{\delta}^{n},
\]
which is (I.16).  The full grid $F^n$ hits every polynomial of
individual degree less than $q$: induction on the number of variables,
using the univariate root bound, shows that a polynomial vanishing
on that grid must be zero.  This also gives the matching upper bound
when $\delta=q-1$.  For $\delta\ge q$, even that grid cannot be a
hitting set for all formally nonzero polynomials, because
$x_1^q-x_1$ vanishes on it. \hfill$\square$

\paragraph{Precise scope.}
The lower bound concerns evaluations in the coefficient field.
It does not apply to extension-field evaluations, which can behave
very differently for a fixed base-field coefficient class.
The row-state lift is useful as an exact finite arithmetic simulation
and avoids the erroneous unweighted support graph of I-02.
It does not construct the graph from an evaluation oracle or make
its state-dependent path weights into variable cocharacters.

\subsection{I-17. A concrete uniform hitting set for invertible factors}
\label{I-17}

\paragraph{Background and statement.}
Invertible-factor ROABP hitting sets are already a result of
Agrawal et al. (2014), Theorem~3, pp.~5 and~19, and their
Lemma~27, p.~22, gives matrix coefficient concentration from
invertible constant terms.  The following is a self-contained
row-space specialization with constant boundary vectors; no
historical-priority claim is made.

Let $n,w\ge1$ and $\delta\ge1$.  Suppose
\[
 f=u^{\mathsf T}A_1(x_{\pi(1)})\cdots A_n(x_{\pi(n)})v
\]
has square $w\times w$ factors of individual degree at most
$\delta$, each with nonzero determinant polynomial.
Choose sets $S,B\subseteq F$ with
$|S|=nw\delta+1$, $|B|=\delta+1$, and $0\in B$.
Then the following single set, independent of the factors and order,
hits every nonzero such $f$:
\[
 \mathcal H=\{a\mathbf1+b: a\in S,\ b\in B^n,
                         |\operatorname{supp}(b)|\le w-1\}.
\tag{I.17}\label{eq:I-invertible-hit}
\]
Its cardinality is at most
\[
 (nw\delta+1)\sum_{j=0}^{\min(w-1,n)}\binom nj\delta^j.
\]
It is polynomial in $n,\delta$ for every fixed $w$, works in arbitrary
characteristic with the stated field-size condition, and requires
no knowledge of $\pi$.

\paragraph{Proof.}
The nonzero polynomial $\prod_i\det A_i(t)$ has degree at most
$nw\delta$, so some $a\in S$ makes every $A_i(a)$ invertible.
Reindex the variables in their true order for this proof only.
Set $P_0=I$, $P_i=A_1(a)\cdots A_i(a)$, and
\[
 M_i(y)=P_{i-1}A_i(a+y)P_i^{-1}.
\]
Then $M_i(0)=I$ and telescoping gives
\[
 f(a\mathbf1+\mathbf y)
    =u^{\mathsf T}M_1(y_1)\cdots M_n(y_n)P_nv.
\]
For nonzero $f$ we have $u\ne0$.  Let $R_i\subseteq F^{1\times w}$
be the span of the coefficient rows of
$u^{\mathsf T}M_1(y_1)\cdots M_i(y_i)$, and put $d_i=\dim R_i$.
Because the next constant term is $I$, these spaces are nested and
$d_0=1$, $d_i\le w$.

We claim that $R_i$ has a basis consisting of actual coefficient
rows whose monomials have support at most $d_i-1$.  This is true for
$i=0$.  Inductively retain a basis for $R_{i-1}$.  If the dimension
does not grow, that basis suffices.  Otherwise extend it using rows
$rM_{i,a}$, where $r$ is a retained coefficient row and $a>0$.
These rows span all new coefficients, and each is itself a coefficient
row with support at most $d_{i-1}$.  Since $d_i>d_{i-1}$, every
retained or added basis row has support at most $d_i-1$.
This proves the claim.

At $i=n$, at least one coefficient row has nonzero product with
$P_nv$.  Therefore some basis coefficient of support at most $w-1$
also has nonzero product.  Restricting the shifted scalar polynomial
to those variables, with the others zero, gives a nonzero polynomial
of individual degree at most $\delta$.  The grid $B$ on those
variables hits it by repeated univariate root counting.  This point
lies in (I.17).  The set (I.17) is invariant under coordinate
permutations, so its construction never needed the true order.
Counting supports proves its size bound. \hfill$\square$

The support bound is sharp for the normalized row-space argument:
for $w\ge2$ and $n=w-1$, factors $I+x_iE_{i,i+1}$ and boundaries $e_1,e_w$
compute $x_1\cdots x_{w-1}$.  For $w=1$, the theorem is the familiar
diagonal univariate test.  Generically singular factors require a
different argument, as the following explicit family shows.

\subsection{I-18. Regularization must preserve a unit Pl\"ucker coordinate}
\label{I-18}

\paragraph{Statement and counterexample.}
Let $m\ge2$, $n=2m$, and
\[
 f(\mathbf x)=\prod_{j=1}^m(x_{2j}-x_{2j-1}).
\]
This is a nonzero multilinear width-two ROABP in the displayed
variable order.  It vanishes at every point $a\mathbf1+b$ with
$|\operatorname{supp}(b)|<m$, for every scalar $a$.
Consequently the sampling restriction in I-17 cannot extend to
all width-two ROABPs by regularizing their matrices and specializing
the regularization parameter to zero.

More precisely, let $R=F[[\varepsilon]]$, $K=F((\varepsilon))$,
$C(\varepsilon)\in R^{r\times s}$ have
$\operatorname{rank}C(0)=r$, and let $P\in R^{s\times h}$ be any
integral compression.  Then
\[
 \operatorname{rank}_K(CP)=r
 \quad\Longleftrightarrow\quad
 \text{some $r$-minor of $CP$ has finite valuation},
\]
whereas
\[
 \operatorname{rank}_F(C(0)P(0))=r
 \quad\Longleftrightarrow\quad
 \text{some $r$-minor of $CP$ has valuation zero}.
\tag{I.18}\label{eq:I-unit-Plucker}
\]
The latter is equivalent to $\operatorname{im}_R(CP)=R^r$.
Generic span equality alone is insufficient.

\paragraph{Proof.}
Use the two singular matrix factors
\[
 L(x)=\begin{pmatrix}-x&1\\0&0\end{pmatrix},\qquad
 R(y)=\begin{pmatrix}1&0\\y&0\end{pmatrix}.
\]
Their product has $(1,1)$ entry $y-x$ and all other entries zero.
Taking the ordered product of these pairs with boundaries $e_1,e_1$
computes $f$.  At $a\mathbf1+b$ each pair difference is
$b_{2j}-b_{2j-1}$.  If fewer than $m$ coordinates of $b$ are
nonzero, one of the $m$ disjoint pairs has both entries zero,
so the product vanishes.

Replace every matrix factor by that factor plus $\varepsilon I$.
All regularized factors are generically invertible over $K$,
because their determinant is monic in $\varepsilon$.
Call the resulting scalar polynomial $f_\varepsilon$.
At the zero point a regularized pair is
\[
 (L(0)+\varepsilon I)(R(0)+\varepsilon I)
  =\begin{pmatrix}\varepsilon(1+\varepsilon)&\varepsilon\\
                   0&\varepsilon^2\end{pmatrix},
\]
so $f_\varepsilon(0)=\varepsilon^m(1+\varepsilon)^m\ne0$ in $K$.
Yet every common-shift sample of perturbation support less than $m$
has zero specialization at $\varepsilon=0$.

For the matrix assertion, a minor is nonzero over $K$ exactly when
it has finite valuation, and its reduction is nonzero over $F$
exactly when its valuation is zero.  This proves the two rank
criteria.  A unit $r$-minor gives $r$ columns invertible over $R$,
so their span is $R^r$.  Conversely, if the columns span $R^r$,
their reductions span $F^r$, which supplies a unit minor.
\hfill$\square$

\paragraph{The resulting special-fiber gate.}
In the scalar example the module generated by all coefficients of
$f_\varepsilon$ is $R$, since some coefficient of $f_0=f$ is nonzero.
For any finite collection of the displayed fixed $F$-valued samples
that contains the zero point, their values generate a nonzero ideal $N\subseteq
\varepsilon R$.  Hence $K N=K$ and its saturation
$(KN)\cap R$ is $R$, although all original sampled values are zero.
Saturating by dividing by a power of $\varepsilon$ thus does not
supply a nonzero evaluation of the original black box.
For vector coefficients the correct next target is a compression
with a unit Pl\"ucker coordinate, as in (I.18), or an equivalent
proof of coefficient-span preservation on the special fiber.
This refines the generic Grassmannian condition in I-14; it is not
established merely by flatness of a surrounding Rees family.

\subsection{Open gate I-G01 and the order of further work}
\label{I-G01}

The surviving main gate is an explicit polynomial family of
integer cocharacters of polynomial magnitude which, for every
width-$w$ matrix product with disjoint univariate factors,
enters at least one cone (I.3), or satisfies the exact weaker joint
coefficient-span condition I-14 for the family.  If this is
proved, I-B02 gives polynomial hitting sets, and I-08 gives the
precisely delimited matrix-invariant conclusion.  No finding
above claims that this gate is closed.

The constructions already attempted and resolved as stated are:
direct state weights (I-02); monomial-alphabet lifting followed
by exact evaluation interpolation (I-06); positive Gram basis
enumeration (I-05); independent and correlated Haar repair
(I-07); exterior/permutation expansion and literal subspace
configuration graphs (I-11); and positive coefficient-norm
transfer matrices (I-09), exact universal energy cubature (I-13),
and the arithmetic support test (I-12).  The further repairs now
resolved are joint span preservation (I-14), bounded affine mask
decoding (I-15), exact finite-field row states (I-16), the restricted
invertible-factor sampler (I-17), and epsilon regularization with
its necessary special-fiber condition (I-18).  The positive result remaining from
direct graph isolation is I-01.  A next substantive construction
would have to preserve the monomial quotient and coefficient
dependences while compressing the matroid-cone search, rather
than merely isolate another term in a cancellation-prone
expansion.
